% !TEX program = xelatex
% =========================================================
%  中文译稿模板 -- pdf-to-chinese-latex skill
%
%  使用方法：
%   1. 把所有 <...> 占位符替换为译稿内容；
%   2. 在 images/ 目录下放置 fig1.png, fig2.png, ... ；
%   3. Overleaf 编译方式：Menu -> Compiler 选 XeLaTeX -> Recompile。
%   4. 中英对照版：同目录放 main_bilingual.tex（skill 的 references/ 里有），
%      编译它即得“上原文下中文”的 main_bilingual.pdf；正文只在本文件维护。
%
%  本模板预加载的包足以支撑大部分学术论文：
%   - ctex          中文支持（XeLaTeX）
%   - amsmath/amssymb/amsthm  公式
%   - graphicx + booktabs + tabularx + longtable + multirow  图与表
%   - algorithm + algpseudocode  算法伪代码
%   - hyperref + url  超链接
%   - enumitem    自定义列表
% =========================================================
\documentclass[a4paper,11pt]{article}

\usepackage[UTF8,heading=true]{ctex}
\usepackage[a4paper,margin=2.5cm]{geometry}
\usepackage{amsmath,amssymb,amsthm}
\usepackage{graphicx}
\usepackage{booktabs}
\usepackage{array}
\usepackage{longtable}
\usepackage{caption}
\usepackage{enumitem}
\usepackage{algorithm}
\usepackage{algpseudocode}
\usepackage{xcolor}
\usepackage[hidelinks,colorlinks=false]{hyperref}
\usepackage{url}
\usepackage{tabularx}
\usepackage{multirow}
\usepackage{float}        % provides the [H] specifier to FORCE a figure
                          % into position (no floating). Use only when a
                          % regular [!htbp] still defers the figure.
\usepackage{flafter}      % a float never lands above the heading / text that precedes it

% 西文里的 en dash（如 "Mistral 7B – Instruct"）按西文排，不当成中文标点吃掉两侧空格（中文破折号 —— 是 U+2014，不受影响）
\xeCJKDeclareCharClass{Default}{"2013}
% 放宽浮动体参数：一页可多放几张图表，避免只放一张图的半空浮动页
\renewcommand{\topfraction}{0.9}
\renewcommand{\bottomfraction}{0.8}
\renewcommand{\textfraction}{0.07}
\renewcommand{\floatpagefraction}{0.8}
\setcounter{topnumber}{3}
\setcounter{totalnumber}{4}
% 中西文混排时允许多一点行间伸缩，减少 overfull 行（长英文术语、URL）
\setlength{\emergencystretch}{2em}

\graphicspath{{images/}}  % so \includegraphics{fig1.png} resolves to images/fig1.png

% --- Figure inclusion convention (read before pasting figures) -----------
% Default to [!htbp] placement and pick width from the PNG's aspect ratio:
%
%   aspect ≥ 2.5  → \includegraphics[width=0.85\textwidth]{figN.png}  (very wide banner)
%   1.5 ≤ < 2.5   → \includegraphics[width=0.75\textwidth]{figN.png}  (standard landscape)
%   0.9 ≤ < 1.5   → \includegraphics[width=0.65\textwidth]{figN.png}  (square-ish)
%   < 0.9         → \includegraphics[width=0.55\textwidth]{figN.png}  (portrait)
%
% render_figures.py emits an `_figure_includes.tex` sidecar with these
% widths already filled in for every extracted figure — paste from there.
% -------------------------------------------------------------------------

\graphicspath{{images/}}
\captionsetup{font=small,labelfont=bf}
\setlength{\parskip}{0.4em}

% --- 中英对照宏（纯中文版定义）-------------------------------------------
% 正文的每个段落 / 标题 / 列表项 / 脚注都同时写原文和译文：
%   \bipar{原文}{中文}  \bisection{原文}{中文}  \bihead{原文}{中文}
%   \biitem{原文}{中文} \bifoot{原文}{中文}
% 编译 main.tex 时下面的定义生效，只排中文（= 普通译文版 main.pdf）。
% 编译 main_bilingual.tex 时，它预先定义了“上原文下中文”的版本，
% 这里的 \providecommand 不会覆盖它（-> main_bilingual.pdf）。
% 所以原文一侧也必须是合法 LaTeX：% & # _ $ 要转义，数学写成 $...$。
\providecommand{\bipar}[2]{#2}
\providecommand{\bihead}[2]{#2}
\providecommand{\bisection}[2]{\section{#2}}
\providecommand{\bisubsection}[2]{\subsection{#2}}
\providecommand{\bisubsubsection}[2]{\subsubsection{#2}}
\providecommand{\biitem}[2]{\item #2}
\providecommand{\bifoot}[2]{\footnote{#2}}
\providecommand{\bithanks}[2]{\thanks{#2}}
% “摘要”“参考文献”标题不用手动改：对照版入口会自动配上 Abstract / References。
% -------------------------------------------------------------------------

% 自定义环境
\newtheorem{definition}{定义}
\newtheorem{assumption}{假设}
\newtheorem{theorem}{定理}
\newtheorem{lemma}{引理}
\newtheorem{proposition}{命题}

\title{\bfseries \bihead{<原文标题>}{<中文标题>}}
% 作者脚注（贡献说明、“工作完成于…”等）用 \bithanks{<原文>}{<中文>}，不要用 \bifoot
\author{<罗马字作者 1>\bithanks{<原文作者脚注>}{<中文作者脚注>}\quad <罗马字作者 2>\\[2pt]
\small\itshape <作者机构 1>\\
\small\itshape <作者机构 2>\\
\small 通讯作者邮箱：\href{mailto:<email>}{<email>}}
\date{原文发表于 \emph{<期刊全名>}，<年份>，第 <卷> 卷 第 <期> 期，pp.\ <起>--<止>\\
\small DOI:\ \href{<DOI URL>}{<DOI>}}

\begin{document}
\maketitle

\begin{abstract}
% \noindent 写在中文那一半里面（写在 \bipar 前面会在对照版里多出一个空行）
\bipar{<原文摘要>}{\noindent <摘要正文翻译>}

\bipar{\textbf{Keywords:} <keyword 1>; <keyword 2>; <keyword 3>}{\noindent\textbf{关键词：}<关键词 1>；<关键词 2>；<关键词 3>}
\end{abstract}

% =========================================================
\bisection{Introduction}{引言}
\label{sec:intro}

% 每个原文段落对应一个 \bipar，段与段之间照常空一行。
% 脚注只写在中文那一半里：\bifoot{<原文脚注>}{<中文脚注>}。
\bipar{<引言第 1 段原文>}{<引言第 1 段译文>\bifoot{<原文脚注>}{<中文脚注>}}

\bipar{<引言第 2 段原文>}{<引言第 2 段译文>}

% =========================================================
\bisection{<原文第二节标题>}{<第二节标题>}
\label{sec:methods}

\bipar{<原文段落>}{<正文翻译...>}

% --- 公式示例 ---
% 用 \label{eq:cN} 对应原文的 (N) 式，正文里用 \eqref{eq:cN} 引用。
% 公式、表格内容、图片、算法、参考文献只写一份，不进 \bipar。
\begin{equation}
\label{eq:c1}
\min z = \sum_{i=1}^{n} c_i x_i.
\end{equation}

% 被公式隔开的段落片段：公式前后各一个 \bipar，公式紧贴前一个片段（不空行）。
\bipar{The constraints are given in Eq.~\eqref{eq:c2}:}{约束如式 \eqref{eq:c2} 所示：}
\begin{equation}
\label{eq:c2}
\sum_{j=1}^{m} a_{ij} x_j \le b_i,\quad \forall i.
\end{equation}

% --- 图示例 ---
% 裁出的 PNG 只含图本身、不含原文图题（图题由 LaTeX 排），所以图题一律写成
% \caption{\bihead{<原文图题>}{<中文图题>}}（原文图题去掉 "Figure N:" 编号）；
% 图里有文字（标签、图例、坐标轴）时，中文一侧以“图中文字：…”结尾逐项给出译名。
\begin{figure}[!htbp]
\centering
\includegraphics[width=0.75\textwidth]{fig1.png}
\caption{\bihead{<Figure 1 original caption>}{<图 1 的中文标题>。图中文字：<Label A> 标签 A；<Label B> 标签 B}}
\label{fig:f1}
\end{figure}

% --- 表示例（用 tabularx + booktabs）---
% 表格内容只写一份（中文表头）；表题用 \bihead 写原文 + 中文。
\begin{table}[ht]
\centering
\caption{\bihead{<Table 1 original caption>}{<表 1 的中文标题>}}
\label{tab:t1}
\begin{tabularx}{0.85\textwidth}{c X r}
\toprule
参数 & 含义 & 数值 \\
\midrule
$\alpha$ & <含义说明> & 1.0 \\
$\beta$  & <含义说明> & 0.5 \\
\bottomrule
\end{tabularx}
\end{table}

% --- 算法示例 ---
\begin{algorithm}[ht]
\caption{\bihead{<Algorithm name>}{<算法名称>}}
\label{alg:a1}
\begin{algorithmic}[1]
\State \textbf{输入：}<输入说明>
\State \textbf{输出：}<输出说明>
\For{$i=1$ to $n$}
  \If{<条件>}
    \State <操作>
  \EndIf
\EndFor
\State \Return <返回值>
\end{algorithmic}
\end{algorithm}

% =========================================================
\bisection{Conclusion}{结论}
\label{sec:conclusion}

\bipar{<原文结论>}{<结论翻译...>}

% 列表：每项用 \biitem 代替 \item
\begin{itemize}
\biitem{<原文要点 1>}{<要点 1 译文>}
\biitem{<原文要点 2>}{<要点 2 译文>}
\end{itemize}

% 不编号的标题用 \section*{\bihead{...}{...}}
\subsection*{\bihead{Acknowledgements}{致谢}}
\bipar{<原文致谢>}{<致谢翻译...>}

% =========================================================
\begin{thebibliography}{99}
\small
\bibitem{Key1} <Author>, <Title>, <Journal>, <Year>. doi: <DOI>.
\bibitem{Key2} <Author>, <Title>, <Journal>, <Year>. doi: <DOI>.
\end{thebibliography}

\end{document}
